\documentclass[11pt]{article}
\usepackage[margin=1in]{geometry}
\usepackage{amsmath}
\usepackage{parskip}
\begin{document}
\thispagestyle{empty}

\noindent Journal of Fluid Mechanics\\
\noindent Attn.\ Professor Sutanu Sarkar, Associate Editor

\bigskip
\noindent Re: revised submission developed from JFM-2026-1781

\bigskip
\noindent Dear Professor Sarkar,

We submit for consideration a revised manuscript, developed from JFM-2026-1781,
which was declined in August 2026. Following the editorial policy stated in that
decision, we cite the original submission's number and ask that the manuscript be
assigned initially to the associate editor who handled it.

The three reports converged on two points. First, the central spectral prediction
had been tested only against the model's own kinematics, so the agreement carried
little weight. Second, the paper was long and did not separate its new results from
established ones. We have rebuilt the manuscript around both.

The model is now measured against an independent benchmark. A $128^3$ strained-box
direct numerical simulation, together with exact rapid-distortion companions reduced
to the same one-dimensional observables, provides the spectral comparison
(\S4.6, figure~11), and the departure from linear theory is reported as a function
of strain rapidity (figure~12). The linear reference is now the exact
rapid-distortion solution projected onto the line, which is broadband rather than a
rigid spectral translation (\S4.2, figure~8); the rigid-translation object of the
earlier version is identified as the model's own kinematics. This addresses the
benchmark concern raised by all three referees and Referee~3's objection to the
linear reference.

Section~5 no longer offers a closed single-line functional. It states the
obstruction exactly, derives sharp bounds on the rapid pressure--strain from line
data alone, measures the axisymmetry assumption rather than imposing it, and reports
the component-ordering deficiency against the direct simulation, localising it in the
closure. Where the single-line representation falls short, the shortfall is now
quantified.

The leading-edge-noise application is carried through (\S7): the computed distorted
inflow spectrum is supplied to Amiet's theory and the predicted level compared with
four published experiments (figures~17--21).

The manuscript has been reorganised and shortened. Implementation detail and the
catalogue of intervention variants are moved to the supplementary material,
established and new results are separated, the notation is fixed once, and the
introduction has been revised to include the near-edge-distortion literature the
referees identified.

A point-by-point response to each referee accompanies this letter. We are grateful
to the referees; their reports directed the revision toward the validation and the
scope the work needed.

\bigskip
\noindent Yours sincerely,\\
\noindent Sparsh Sharma and Alan R.\ Kerstein

\end{document}
